\section{Introduction}
Machine learning systems are capable of exhibiting discriminatory behaviors against certain demographic groups in high-stakes domains such as law, finance, and medicine \citep{kirchner2016machine,aleo2008foreclosure,kim2015sex}.
These outcomes are potentially unethical or illegal \citep{barocas2014datas,hellman2017indirect}, and behoove researchers to investigate more equitable and robust models.
One promising approach is fair representation learning: the design of neural networks using learning objectives that satisfy certain fairness or parity constraints in their outputs \citep{zemel2013learning,louizos2015variational,edwards2015censoring,madras2018learning}.
This is attractive because neural network representations often generalize to tasks that are unspecified at train time, which implies that a properly specified fair network can act as a group parity bottleneck that reduces discrimination in unknown downstream tasks.

Current approaches to fair representation learning are flexible with respect to downstream tasks but inflexible with respect to sensitive attributes.
While a single learned representation can adapt to the prediction of different task labels~$y$, the single sensitive attribute $a$ for \emph{all} tasks must be specified at train time.
Mis-specified or overly constraining train-time sensitive attributes could negatively affect performance on downstream prediction tasks.
Can we instead learn a \textit{flexibly fair} representation that can be \emph{adapted}, at test time, to be fair to a variety of protected groups and their intersections?
Such a representation should satisfy two criteria.
Firstly, the structure of the latent code should facilitate \textit{simple} adaptation, allowing a practitioner to easily adapt the representation to a variety of fair classification settings, where each task may have a different task label $y$ and sensitive attributes $a$.
Secondly, the adaptations should be \textit{compositional}: the representations can be made fair with respect to conjunctions of sensitive attributes, to guard against subgroup discrimination (e.g., a classifier that is fair to women but not Black women over the age of 60).
This type of subgroup discrimination has been observed in commercial machine learning systems \citep{buolamwini2018gender}.

In this work, we investigate how to learn flexibly fair representations that can be easily adapted at test time to achieve fairness with respect to sets of sensitive groups or subgroups.
We draw inspiration from the disentangled representation literature, where the goal is for each dimension of the representation (also called the ``latent code'') to correspond to no more than one semantic factor of variation in the data (for example, independent visual features like object shape and position) \citep{higgins2016beta, locatello2018challenging}.
Our method uses multiple sensitive attribute labels at train time to induce a disentangled structure in the learned representation, which allows us to easily eliminate their influence at test time.
Importantly, at test time our method does not require access to the sensitive attributes, which can be difficult to collect in practice due to legal restrictions \citep{elliot2008new,division1983}.
The trained representation permits simple and composable modifications at test time that eliminate the influence of sensitive attributes, enabling a wide variety of downstream tasks.

We first provide proof-of-concept by generating a variant of the synthetic DSprites dataset with correlated ground truth factors of variation, which is better suited to fairness questions.
We demonstrate that even in the correlated setting, our method is capable of disentangling the effect of several sensitive attributes from data, and that this disentanglement is useful for fair classification tasks downstream.
We then apply our method to a real-world tabular dataset (Communities \& Crime) and an image dataset (Celeb-A), where we find that our method matches or exceeds the fairness-accuracy tradeoff of existing disentangled representation learning approaches on a majority of the evaluated subgroups.
